\item \textbf{Fusión limpia de protón con Boro-11.}
Una reacción de fusión sumamente atractiva para la generación de energía consiste en la absorción de un protón por un núcleo de boro-11 para producir tres partículas alfa sin liberar neutrones radiactivos:
$$ ^{1}_{1}\text{H} + ^{11}_{5}\text{B} \to 3(^{4}_{2}\text{He}) $$
\begin{enumerate}
    \item ¿Cuánta energía en $\text{MeV}$ se libera en cada evento de fusión individual?
    \item ¿Por qué las partículas reactivas deben poseer energías cinéticas extremadamente altas para que ocurra la fusión, a diferencia de la fisión inducida por neutrones lentos?
\end{enumerate}